抛一枚公平硬币 100 次，正面数 \(X\sim\Binomial(100,0.5)\)。大数定律说 \(\frac{X}{100}\to 0.5\)（概率意义下），但实际的 \(\frac{X}{100}\) 可能等于 0.47。中心极限定理进一步说，标准化误差 \(\frac{X-50}{\sqrt{25}}=2(\frac{X}{100}-0.5)\cdot 10\) 近似于 \(N(0,1)\)。例如 \(P(|\frac{X}{100}-0.5|<0.05)=P(|X-50|<5)\approx\Phi(1)-\Phi(-1)\approx 0.68\)——大数定律断言收敛，中心极限定理给出收敛后残余波动的尺度。

单次实验不可预测，并不意味着大量重复没有结构。样本均值会靠近总体均值，标准化误差常呈 Gaussian，模拟平均也能逼近难算积分。但``靠近''有多种含义，``大量''也不自动修复偏差、依赖或无限方差。

本单元依次回答四个问题：

\begin{enumerate}
\tightlist
\item
  随机变量的多种收敛：同一概率空间上的逐路径接近、坏事件概率消失、均方误差消失和分布形状接近有何区别。
\item
  大数定律与长期平均：为什么独立重复样本的平均会稳定，以及弱、强结论分别给出什么。
\item
  中心极限定理与 Gaussian 近似：平均虽趋于常数，它周围剩余的 \(1/\sqrt n\) 级波动长什么样。
\item
  Monte Carlo：从随机样本到数值答案：怎样把大数定律和中心极限定理变成可报告误差的计算方法。
\end{enumerate}

大数定律回答一致性，中心极限定理回答误差尺度与近似形状。二者作用不同，不可混为一谈。

\subsection{一个具体算例}

\(X_i\sim\iid\Uniform(0,1)\)，\(\E[X_i]=0.5\)，\(\Var(X_i)=1/12\)。\(n=100\) 时，WLLN：\(\bar X_{100}\approx 0.5\)。CLT：\(\sqrt{100}(\bar X_{100}-0.5)/\sqrt{1/12}\approx N(0,1)\)，即 \(\bar X_{100}\approx N(0.5,\frac{1}{1200})\)。95\% 置信区间：\(0.5\pm1.96\cdot\sqrt{1/1200}\approx 0.5\pm0.057\)。

Monte Carlo 例子：估计 \(\pi\)。向 \([-1,1]^2\) 随机投点 \(n=10^6\) 次，落在单位圆内的比例 \(\hat p\) 满足 \(\hat\pi=4\hat p\)。标准误差约 \(\sqrt{\pi(4-\pi)/n}\approx 0.0016\)，故 \(\hat\pi\) 的 95\% 区间宽度约 \(\pm 0.003\)。

\subsection{怎么读这一章}

首次阅读以前三篇为主：先区分几乎处处、依概率、\(L^p\) 与依分布收敛，再分清大数定律的``平均靠近目标''和中心极限定理的``标准化误差趋近某种形状''。特征函数证明适合第二遍，用来理解独立性和有限二阶矩分别进入哪一步。Monte Carlo 是应用闭环，必须连同抽样目标、方差、相关性和误差报告一起读；样本量增大只能压缩随机误差，不能修复偏差或错误目标。

\subsection{本章篇目}

以下篇目按概念依赖排列；若从具体问题进入，也可以先打开最接近当前困惑的一篇，再沿篇内链接回补。

\begin{itemize}
\tightlist
\item \hyperref[sec:05-Probability/08-Limit-Theorems/01]{``随机变量的多种收敛''}；
\item \hyperref[sec:05-Probability/08-Limit-Theorems/02]{``大数定律与长期平均''}；
\item \hyperref[sec:05-Probability/08-Limit-Theorems/03]{``中心极限定理与 Gaussian 近似''}；
\item \hyperref[sec:05-Probability/08-Limit-Theorems/04]{``特征函数怎样证明中心极限定理''}；
\item \hyperref[sec:05-Probability/08-Limit-Theorems/05]{``Monte Carlo：从随机样本到数值答案''}。
\end{itemize}
